\documentclass[10pt,a4paper]{amsart}
\usepackage[margin=19mm]{geometry}
\usepackage{amsmath,amssymb,mathtools}
\usepackage[scr=rsfs,cal=euler]{mathalfa}
\usepackage{xfrac}
\usepackage[unicode,hidelinks,bookmarks=false]{hyperref}
\newcommand{\ie}{i.e.\ }
\newcommand{\cf}{cf.\ }
\newcommand{\resp}{respectively\ }
\newcommand{\Subsection}{\subsection}
\providecommand{\nobreakdash}{\nobreak\hbox{-}}
\setlength{\emergencystretch}{2em}
\multlinegap0pt
\newtheorem{proposition}{Proposition}
\newcommand{\Margin}[1]{}
\def\Mpp{M_{\perp\!\!\!\!\!\bigcirc}}
\def\Mpptil{\tilde{M}_{\perp\!\!\!\!\!\bigcirc}}
\renewcommand\Re{\mathrm{Re}}
\def\fpp{f_{\perp\!\!\!\!\!\bigcirc}}
\def\pp{p_{\perp\!\!\!\!\!\bigcirc}}
\def\pptil{\tilde{p}_{\perp\!\!\!\!\!\bigcirc}}
\newcommand{\revision}[1]{{ #1}}
\newcommand{\skipsimple}[1]{}
\newcommand{\uld}[1]{\underline{d#1}}
\def\uldMpp{\uld{M}_{\perp\!\!\!\!\!\bigcirc}}
\def\ulddMpp{\uld{dM}_{\perp\!\!\!\!\!\bigcirc}}
\title{On uniqueness of coefficient identification in the Bloch-Torrey equation for magnetic resonance imaging}
\author{Barbara Kaltenbacher}
\date{Source: arXiv:2506.13708v2 (CC BY 4.0)}
\begin{document}
\maketitle

\noindent\emph{Source and licence.} On uniqueness of coefficient identification in the Bloch-Torrey equation for magnetic resonance imaging. Version arXiv:2506.13708v2 (CC BY 4.0), distributed by arXiv under the Creative Commons Attribution 4.0 International licence, http://creativecommons.org/licenses/by/4.0/, as recorded in the arXiv licence metadata for this exact version and shown on the abstract page https://arxiv.org/abs/2506.13708v2. Full licence notice: \texttt{licenses/arxiv-2506.13708-CC-BY-4.0.txt}.\par\smallskip

\noindent\textit{Editorial scope.} Counted unit: the article's proposition prop:perturb\_gen\_deriv (source0.tex lines 1365--1375). The article's complete original proof of it is delivered with it (source0.tex lines 1376--1409, one \texttt{proof} environment of 34 lines). No mathematics is written, repaired or re-derived: the statement and the proof are the article's own lines, in the article's own notation, and no retained line is substituted or re-flowed.  Retained uncounted as the source-local context the proof needs: lines 470--481; lines 502--504; lines 508--517; lines 520--532; lines 534--541; lines 547--591; lines 1263--1269; lines 1270--1280; lines 1281--1286.  Nothing was omitted from any retained span, so the package carries no omission record.  No retained line contains a citation, so the wrapper carries no bibliography and no bibliography entry.  No retained line contains a figure environment, an image-inclusion command or drawing code, so no image is delivered and the package declares no asset dependency.  Editorial formatting only, in addition: an AMS \texttt{amsart} preamble that restates the article's own declarations and macros used by the retained lines, namely the environments \texttt{proposition} on separate counters, together with the standard packages those lines need.  The retained environments print in this document as Proposition 1 in order of appearance, whereas the article prints the same items with the numbering of its own full document.  The complete native source of the article as distributed in the arXiv e-print for this exact version is bundled unchanged as \texttt{sources/179966/source0.tex}.
\begin{equation}\label{eqn:bloch-torrey-complex_f}
\begin{aligned}
&\frac{d}{dt}M_\perp(t, \vec r)+{\vec v}(t, \vec r)\cdot\nabla M_\perp(t, \vec r) -\nabla\cdot\bigl(D(\vec r) \nabla M_\perp(t, \vec r)\bigr) \\
&+\bigl(R_2(\vec r)+i(\omega_0 + \gamma\delta B^0(\vec r) + \gamma\vec r \cdot \vec G(t))
\bigr)M_\perp(t, \vec r)-i\gamma M_z(t, \vec r)c^+(\vec r)p(t)=f_\perp(t, \vec r)\\    
&\frac{d}{dt}M_z(t, \vec r)+{\vec v}(t, \vec r)\cdot\nabla M_z(t, \vec r) -\nabla\cdot\bigl(D(\vec r) \nabla M_z(t, \vec r)\bigr) \\
&+R_1(\vec r) M_z(t, \vec r)
-\gamma\Re\bigl(iM_\perp(t, \vec r)\overline{c^+(\vec r)p(t)}\bigr)=f_z(t, \vec r)
\\
&t>0, \quad \vec r\in\Omega,
\end{aligned}
\end{equation}

\begin{equation}\label{divv}
{\vec v}={\vec v}_0+\delta {\vec v}\text{ such that }\nabla\cdot{\vec v}_0\leq0\text{ a.e. in }(0,T)\times\Omega, \quad \|\nabla\cdot\delta{\vec v}\|_{L^1(0,T;L^\infty(\Omega))}<1. 
\end{equation}

\begin{equation}\label{eqn:bloch-torrey-complex_rot_f}
\begin{aligned}
&\frac{d}{dt}\Mpp(t, \vec r)+{\vec v}(t, \vec r)\cdot\nabla \Mpp(t, \vec r) -\nabla\cdot\bigl(D(\vec r) \nabla \Mpp(t, \vec r)\bigr) \\
&+\bigl(R_2^*(\vec r)+i\gamma\vec r \cdot \vec G(t)\bigr)\Mpp(t, \vec r)
-i\gamma M_z(t, \vec r)c^+(\vec r)\pp(t)= \fpp(t, \vec r)\\    
&\frac{d}{dt}M_z(t, \vec r)+{\vec v}(t, \vec r)\cdot\nabla M_z(t, \vec r) -\nabla\cdot\bigl(D(\vec r) \nabla M_z(t, \vec r)\bigr) \\
&+R_1(\vec r) M_z(t, \vec r)
-\gamma\Re\bigl(i\Mpp(t, \vec r)\overline{c^+(\vec r)\pp(t)}\bigr)=f_z(t, \vec r)
\end{aligned}
\end{equation}

\begin{equation}\label{HD1}
\begin{aligned}
&\langle (v_1,w_1),(v_2,w_2)\rangle_{H_D^1} :=
\int_\Omega\Bigl\{ \Re\Bigl((D(\vec r) \nabla v_1(\vec r))\cdot \overline{\nabla v_2(\vec r)}
+R_2(\vec r)\,v_1(\vec r))\, \overline{v_2(\vec r)}\Bigr)\\    
&\phantom{\langle (v_1,w_1),(v_2,w_2)\rangle := \int_\Omega\Bigl(}
+(D(\vec r) \nabla w_1(\vec r))\cdot \nabla w_2(\vec r)
+R_1(\vec r) w_1(\vec r)w_2(\vec r)\Bigr\}\, d \vec r\\
&\phantom{\langle (v_1,w_1),(v_2,w_2)\rangle :=}
+\int_{\Gamma_I}\Bigl\{\Re\Bigl(\beta_\perp(\vec r)\,v_1(\vec r))\, \overline{v_2(\vec r)}\Bigr)
+\beta_z(\vec r)\,w_1(\vec r))\, w_2(\vec r)\Bigr\}\, dS(\vec r).
\end{aligned}
\end{equation}

\begin{equation}\label{definiteness}
\begin{aligned}
&D(\vec r)\text{ nonnegative definite }, \quad 
R_1(\vec r)\geq \underline{R}_1>0, \quad
R_2(\vec r)=\Re(R_2^*(\vec r))\geq \underline{R}_2>0\\
&\vec r\in\Omega,
\end{aligned}
\end{equation}

\begin{proposition}\label{prop:wellposed}
%Let $\omega_0$, $\gamma\in\mathbb{R}$, and 
%$R_1,\, R_2,\,\delta B^0\in L^\infty(\Omega;\mathbb{R})$, (that is, $R_2^*\in L^\infty(\Omega;\mathbb{C})$, 
%$\vec G\in L^1(0,T;\mathbb{R}^3))$, 
%$\vec v\in L^1(0,T;W^{1,\infty}(\Omega))$ satisfying \eqref{divv},
%$D\in L^\infty(\Omega;\mathbb{R}^{3\times 3})$ with \eqref{definiteness},
%$p\in L^1(0,T)$ 
Let $T\in(0,\infty]$, $\gamma\in\mathbb{R}$, 
$R_1\in L^\infty(\Omega;\mathbb{R})$, $R_2^*,\, c^+\in L^\infty(\Omega;\mathbb{C})$, 
$\vec v\in L^1(0,T;W^{1,\infty}(\Omega;\mathbb{R}^3))$, 
$D\in L^\infty(\Omega;\mathbb{R}^{3\times 3})$, 
satisfying \eqref{divv}, \eqref{definiteness},
$\vec G\in L^1(0,T;\mathbb{R}^3)$, 
$\pp\in L^1(0,T;\mathbb{C})$, 
and 
\Margin{R2 6)}
$\revision{\vec{f}=(f_\perp,f_z)}\in L^1(0,T;L^2(\Omega;\mathbb{C}\times\mathbb{R}))+L^2(0,T;H_D^1(\Omega;\mathbb{C}\times\mathbb{R})^*)$.\footnote{Here ${}^*$ denotes the topological dual, for the definition of the space $H_D^1(\Omega)$ we point to \eqref{HD1}, and $f\in W_1+W_2$ means $f=f_1+f_2$ for some $f_1\in W_1$, $f_2\in W_2$.} 

Then the system 
%\eqref{eqn:bloch-torrey-complex_f}
\eqref{eqn:bloch-torrey-complex_rot_f}
has a unique weak solution
%\footnote{with the time derivative moved to the test function by integration by parts}
\begin{equation}\label{V_wellposed}
(\Mpp,M_z)\in \mathbb{V}:=L^\infty(0,T;L^2(\Omega;\mathbb{C}\times\mathbb{R}))\cap L^2(0,T;H_D^1(\Omega;\mathbb{C}\times\mathbb{R}))
\end{equation}
and the estimate 
\begin{equation}\label{enest_thm}
\begin{aligned}
&\|(\Mpp,M_z)\|_{L^q(0,T;L^2(\Omega))}
+\|(\Mpp,M_z)\|_{L^2(0,T;H_D^1(\Omega))}\\
&\leq C\Bigl(\|(\Mpp,M_z)(0, \cdot)\|_{L^2(\Omega)} + \min\{\|\revision{\vec{f}}\|_{L^1(0,T;L^2(\Omega))}, \,  \|\revision{\vec{f}}\|_{L^2(0,T;H_D^1(\Omega;\mathbb{C}\times\mathbb{R})^*)}\}\Bigr)
\end{aligned}
\end{equation}
holds for any $q\in[2,\infty]$ with a constant 
$C$ depending only on $q$ and $1-\|\nabla\cdot\delta{\vec v}\|_{L^1(0,T;L^\infty(\Omega))}$,
but not on $T$ nor the data $(\Mpp,M_z)(0, \cdot)$, $\revision{\vec{f}}$.

If 
$\revision{\vec{f}}\in L^2(0,T;H_D^1(\Omega;\mathbb{C}\times\mathbb{R})^*)$, 
$D^{-1/2}{\vec v}\in L^\infty(0,T;L^\infty(\Omega;\mathbb{R}^{3\times 3})$, 
$\vec G\in L^2(0,T;\mathbb{R}^3)$,
$\pp\in L^2(0,T;\mathbb{C})$, 
then additionally $(\Mpp,M_z)\in H^1(0,T;H_D^1(\Omega;\mathbb{C}\times\mathbb{R})^*)$ holds.
\end{proposition}

\begin{equation}\label{estv}
\begin{aligned}
&\|({\vec v}-\vec{\tilde{v}})\cdot\nabla m\|_{L^q(0,T;L^2(\Omega))}
=\|\bigl(D^{-1/2}({\vec v}-\vec{\tilde{v}})\bigr)\cdot  D^{1/2}\nabla m\|_{L^q(0,T;L^2(\Omega))}\\
&\leq \|D^{-1/2}({\vec v}-\vec{\tilde{v}})\|_{L^{\hat{q}}(0,T;L^\infty(\Omega))} \|m\|_{L^2(0,T;H_D^1(\Omega))}, \quad q\in\{1,2\}, \quad\hat{q}= \frac{2q}{2-q}
\end{aligned}
\end{equation}

\begin{equation*}%\label{estD}
\begin{aligned}
&\|-\nabla\cdot\bigl((D-\tilde{D})\nabla m\bigr)\|_{L^2(0,T;H_D^1(\Omega)^*)}\\
&\skipsimple{
=\Bigl(\int_0^T\Bigl(\sup_{w\in C_c^\infty(\Omega)\setminus\{0\}} \frac{1}{\|w\|_{H_D^1(\Omega)}}
\int_\Omega (D^{-1/2}(D-\tilde{D})D^{-1/2}\, D^{1/2}\nabla m) \cdot D^{-1/2}\nabla w\, d{\vec r}\Bigr)^2 \, dt\Bigr)^{1/2}
}
%\\ &
\leq \|D^{-1/2}(D-\tilde{D})D^{-1/2}\|_{L^\infty(\Omega;\mathbb{R}^{3\times 3})}\, \|m\|_{L^2(0,T;H_D^1(\Omega))}
\end{aligned}
\end{equation*}

\begin{equation}\label{estR}
\begin{aligned}
&\|R\,m\|_{L^q(0,T;L^2(\Omega)}\leq \|R\|_{L^q(0,T;L^\infty(\Omega)}\|m\|_{L^\infty(0,T;L^2(\Omega)}, \quad q\in\{1,2\}\\ 
&R\in\{R_1-\tilde{R}_1,\,R_2^*-\tilde{R}_2^*,\,\vec r \cdot (\vec G(t)-\vec{\tilde{G}}(t)),\,(c^+-\tilde{c}^+)\pp+\tilde{c}^+(\pp-\pptil)\}. 
\end{aligned}
\end{equation}

\begin{proposition}\label{prop:perturb_gen_deriv}
Let $\Omega$ to be a bounded Lipschitz domain.
For all ${x}^{img}$, $\tilde{x}^{img}$ $\in\mathbb{X}^{img}$, ${x}^{mod}$, $\tilde{x}^{mod}$ $\in\mathbb{X}^{mod}$, $\mathcal{S}$ is Fr\'{e}chet differentiable at $(x^{img};x^{mod})$ and $(\tilde{x}^{img};\tilde{x}^{mod})$, and there exists $C_{\mathcal{S}'}>0$ depending only on $\|\tilde{x}^{img}\|_{\mathbb{X}^{img}}$, $\|{x}^{mod}\|_{\mathbb{X}^{mod}}$, $\|\tilde{x}^{mod}\|_{\mathbb{X}^{mod}}$ such that 
\[
\begin{aligned}
&\|\mathcal{S}'(x^{img};x^{mod})-\mathcal{S}'(\tilde{x}^{img};\tilde{x}^{mod})\|_{L(\mathbb{X}^{img}\times\mathbb{X}^{mod},\mathbb{V}_\perp)}\\
&\leq C_{\mathcal{S}'}(\|x^{img}-\tilde{x}^{img}\|_{\mathbb{X}^{img}}+\|x^{mod}-\tilde{x}^{mod}\|_{\mathbb{X}^{mod}})
\end{aligned}
\]
holds.
\end{proposition}

\begin{proof}
Fr\'{e}chet differentiability is straightforward to see, relying on the bilinearity of most of the terms and the fact that one of the spaces underlying the only trilinear term (the one involving ${c}^+\pp$) is a Banach algebra (namely $L^\infty(\Omega)$ for ${c}^+$). 

To show the Lipschitz estimate, we use Proposition~\ref{prop:wellposed} together with the fact that
$(\ulddMpp,\uld{dM}_z)=[\mathcal{S}'(x^{img};x^{mod})-\mathcal{S}'(\tilde{x}^{img};\tilde{x}^{mod})](\uld{x}^{img};\uld{x}^{mod})$ solves \eqref{eqn:bloch-torrey-complex_rot_f} with $(\Mpp,M_z)$ replaced by $(\ulddMpp,\uld{dM}_z)$ and 
\[
\begin{aligned}
-\fpp(t, \vec r)=& 
[{\vec v}-\vec{\tilde{v}}](t, \vec r)\cdot\nabla \uldMpp(t, \vec r) 
-\nabla\cdot\bigl([D-\tilde{D}](\vec r) \nabla \uldMpp(t, \vec r)\bigr) \\
&+\bigl([R_2^*-\tilde{R}_2^*](\vec r)+i\gamma\vec r \cdot [{\vec G}-\vec{\tilde{G}}](t)\bigr)\uldMpp(t, \vec r)\\
&-i\gamma \uld{M}_z(t, \vec r)([c^+-\tilde{c}^+](\vec r)\pp(t)+c^+(\vec r)[\pp-\pptil](t))\\
&+\uld{\vec v}(t, \vec r)\cdot\nabla [\Mpp-\Mpptil](t, \vec r) 
-\nabla\cdot\bigl(\uld{D}(\vec r) \nabla [\Mpp-\Mpptil](t, \vec r)\bigr) \\
&+\bigl(\uld{R}_2^*(\vec r)+i\gamma\vec r \cdot \uld{\vec G}(t)\bigr)[\Mpp-\Mpptil](t, \vec r)\\
&-i\gamma [M_z-\tilde{M}_z](t, \vec r)(\uld{c}^+(\vec r)\pp(t)+c^+(\vec r)\uld{\pp}(t))
\\[1ex]    
-f_z(t, \vec r)&
=[{\vec v}-\vec{\tilde{v}}](t, \vec r)\cdot\nabla \uld{M}_z(t, \vec r) 
-\nabla\cdot\bigl([D-\tilde{D}](\vec r) \nabla \uld{M}_z(t, \vec r)\bigr) \\
&+[R_1-\tilde{R}_1](\vec r) (\uld{M}_z(t, \vec r)-M^{eq}(\vec r))\\
&-\gamma\Re\bigl(i\uldMpp(t, \vec r)\overline{([c^+-\tilde{c}^+](\vec r)\pp(t)+c^+(\vec r)[\pp-\pptil](t))}\bigr)\\
&+\uld{\vec v}(t, \vec r)\cdot\nabla M_z(t, \vec r) -\nabla\cdot\bigl(\uld{D}(\vec r) \nabla M_z(t, \vec r)\bigr) \\
&+\uld{R}_1(\vec r) [M_z-\tilde{M}_z](t, \vec r)\\
&-\gamma\Re\bigl(i\Mpp(t, \vec r)\overline{(\uld{c}^+(\vec r)\pp(t)+c^+(\vec r)\uld{\pp}(t))}\bigr),
\end{aligned}
\]
where $(\uldMpp,\uld{M}_z):=\mathcal{S}'(x^{img};x^{mod})(\uld{x}^{img};\uld{x}^{mod})$,
$(\Mpp,M_z):=\mathcal{S}(x^{img};x^{mod})$,
$(\Mpptil,\tilde{M}_z):=\mathcal{S}(\tilde{x}^{img};\tilde{x}^{mod})$.
Each of these terms can be estimated analogously to \eqref{estv}
%, \eqref{estD}, 
-- \eqref{estR}.
\end{proof}

\end{document}
